\vspace*{4cm}
\begin{center}
    {\Huge\bfseries Abstract}
\end{center}
\vskip 1cm

\phantomsection
\addcontentsline{toc}{chapter}{Abstract}

The digital transformation of healthcare has led to the emergence of connected medical devices, electronic health records, remote monitoring platforms, and artificial intelligence-based decision support systems. Although these technologies generate large amounts of heterogeneous data, many eHealth systems still remain fragmented, static, and weakly personalized. Digital Twin technology offers a promising paradigm to address these limitations by creating dynamic virtual representations of physical entities, processes, or patients, continuously updated through data exchange and supported by simulation, analytics, and decision services.

This Master thesis presents a state-of-the-art study of Digital Twins for eHealth. It first introduces Digital Twin fundamentals, including the distinction between digital models, digital shadows, Digital Twins, and intelligent Digital Twins, before presenting the eHealth and digital-health environment in which these principles can be applied. It then analyzes healthcare-oriented Digital Twin applications, architectures, enabling technologies, and implementation challenges, with an expanded discussion of application domains such as personalized medicine, remote monitoring, hospital workflow optimization, public health, and cardiovascular Digital Twins. Particular attention is given to interoperability, data governance, privacy, security, explainability, validation, and ethical considerations. The literature synthesis derives a general multi-layer eHealth Digital Twin architecture, together with data-flow, cross-cutting requirement, and evaluation findings. The final chapter concludes the Master study and introduces CardioTwin as a practical cardiovascular continuation for the PFE, while the thesis remains centered on Digital Twins for eHealth in a general healthcare context.

\textbf{Keywords:} Digital Twin, eHealth, Digital Health, Healthcare, IoT, IoMT, Artificial Intelligence, Simulation, Interoperability, Clinical Decision Support, Cardiovascular Digital Twin.

\clearpage
